\documentclass[11pt]{article}
\usepackage[margin=2.1cm]{geometry}
\usepackage{booktabs}
\usepackage{amsmath}
\usepackage{enumitem}
\usepackage{xcolor}
\usepackage[T1]{fontenc}
\usepackage{lmodern}

\setlength{\parindent}{0pt}
\setlength{\parskip}{3.6pt}
\newcommand{\code}[1]{\texttt{#1}}
\newcommand{\hr}{\vspace{2pt}\hrule\vspace{6pt}}

\begin{document}

\begin{center}
{\Large\bfseries EE\,362 --- Computer Organization Laboratory}\\[2pt]
{\large\bfseries Lab 1 --- Design Flow and the Transparent Latch}\\[4pt]
Fall 2026 \quad$\cdot$\quad Lab TA: Gandhi Torres \quad$\cdot$\quad Sections: Mon / Tue, 3:45--5:45~PM, ECIII 237
\end{center}
\hr

\section*{1. Objective}

Run the complete Quartus and Questa design flow without help, and build the first
circuit in this course that can \emph{remember} something.

By the end of this session you will be able to state, and defend, one sentence:

\begin{quote}
\emph{An equation whose output appears on its own right-hand side is memory.}
\end{quote}

\section*{2. Background: what makes a latch different}

Every circuit in Lab 0 was \textbf{combinational}: the output depended only on the
inputs present at that moment. A transparent latch is different. When its enable
is high it behaves like a wire --- it is \emph{transparent}. When the enable goes
low it \emph{holds} whatever value it had, and stops listening to its data input.

That holding behaviour is memory, and in Boolean terms it shows up as the output
appearing on both sides of its own equation.

\subsection*{Specification --- Table T (transition table)}

\begin{center}
\begin{tabular}{clll@{\hskip 8pt}||@{\hskip 8pt}ll}
\toprule
\# & \code{Clr} & \code{En} & \code{D} & $Q^{+}$ & note \\
\midrule
0 & 1 & -- & -- & \textbf{0} & asynchronous clear, dominates everything \\
1 & 0 & 0 & -- & \textbf{Q} & hold \quad$\leftarrow$\; \emph{the memory row} \\
2 & 0 & 1 & 0 & \textbf{0} & transparent \\
3 & 0 & 1 & 1 & \textbf{1} & transparent \\
\bottomrule
\end{tabular}
\end{center}

\code{--} means \emph{don't care}. \; $Q^{+}$ means \emph{the next value of Q}. \;
A \textbf{Q} in the output column means \emph{hold the value you already had}.

\subsection*{Equations}

You must reproduce this derivation, in three lines, in your report:

\begin{align*}
\text{(1) canonical} \quad & Q^{+} = \overline{Clr}\cdot En\cdot D \;+\; \overline{Clr}\cdot\overline{En}\cdot Q\\
\text{(2) factored} \quad & Q^{+} = \overline{Clr}\cdot\left( En\cdot D + \overline{En}\cdot Q \right) \qquad \text{[distributive law]}\\
\text{(3) VHDL} \quad & \text{the \code{if} statement you are about to write}
\end{align*}

Look at the $Q$ on the right-hand side of line (1). Nothing you wrote in Lab 0 had
that. It is the entire difference between combinational logic and memory.

\textbf{Asynchronous} means \code{Clr} takes effect the instant it is asserted. It
does not wait for \code{En} and it does not wait for a clock --- there is no clock
in this circuit. That is why \code{Clr} is tested \emph{first}: in an
\code{if\,/\,elsif} chain, the first true condition wins.

\section*{3. Guided examples (we do these together)}

\begin{description}[leftmargin=2.4em,style=nextline,itemsep=2pt]
\item[Example 1 --- \code{majority3}] A circuit you already trust, from Lab 0. You
will build its Quartus project \emph{by hand} and walk the entire flow. Never learn
a tool and a circuit at the same time: when something breaks you will not know
which one is wrong.
\item[Example 2 --- \code{trans\_latch}] The transparent latch without the clear
input, plus the testbench pattern used for the rest of the semester.
\item[Example 3 --- \code{trans\_latch\_bad}] The same latch with one signal
missing from its sensitivity list. It \emph{simulates} wrong and \emph{synthesises}
right. Watch this one carefully --- it is the most expensive class of bug in HDL
design, and it costs one missing word.
\end{description}

\section*{4. Assignment}

Start from the files in \code{assignment/}. Open \code{Lab1.qpf} in Quartus.

\begin{enumerate}[leftmargin=2em,itemsep=3pt]
\item \textbf{\code{TransLatch}} --- complete the skeleton. A 1-bit transparent
latch with asynchronous clear, implementing Table T exactly.

\item \textbf{\code{TransLatch8}} --- the same behaviour, 8 bits wide. The control
signals \code{Clr} and \code{En} stay \emph{one bit each} and are shared by all
eight storage elements. Note the pattern: \emph{control signals are narrow, data
paths are wide}. That is true for the rest of this course.

\item \textbf{Testbench} --- complete \code{tb\_TransLatch.vhd}. It already
compiles and runs as supplied; it simply does not check enough yet.
\textbf{At least six checked vectors}, and they must include all of:
\begin{itemize}[leftmargin=1.4em,itemsep=1pt]
\item the \textbf{hold vector}: with \code{En\,=\,'0'}, change \code{D} and prove
\code{Q} does \emph{not} follow. \emph{Without this vector you have not tested that
your design is memory at all --- a plain piece of wire passes every other test you
could write.}
\item clear asserted while the latch is transparent;
\item clear asserted while the latch is holding;
\item at least one \code{TransLatch8} vector using a pattern that is neither
all-zeros nor all-ones (for example \code{x"A5"}), so that a stuck bit is visible.
\end{itemize}
\end{enumerate}

\textbf{Note on the two sensitivity lists in the skeletons.} They are incomplete on
purpose. The files compile as given, and simulate incorrectly --- exactly like
Example 3. Fixing them is part of the assignment.

\section*{5. Evidence to collect (E1, E2, E7)}

See the \emph{Reading Your Results} handout for how to interpret each one.

\begin{center}
\begin{tabular}{lll}
\toprule
& Source & What it must show \\
\midrule
E1 & Flow Summary + Messages & 0 errors; every warning sorted and explained \\
E2 & RTL Viewer & a latch, and nothing you did not intend \\
E7 & Questa waveform + transcript & the hold interval, and 0 assertion failures \\
\bottomrule
\end{tabular}
\end{center}

\section*{6. Report --- \code{lab01.tex} in your Overleaf project}

All seven sections. Specifically for this lab:

\begin{itemize}[leftmargin=1.6em,itemsep=1pt]
\item \textbf{Specification:} Table T, reproduced, with numbered rows.
\item \textbf{Design:} the three-line derivation, with the law named on line (2).
\item \textbf{Verification:} your vector table, and one sentence justifying why
that set of vectors is sufficient.
\item \textbf{Results:} RTL Viewer figure and the waveform, with the hold interval
annotated. \emph{Every figure needs one sentence saying what it proves.} A figure
with no sentence earns zero.
\item \textbf{Discussion:} one bug you hit and how you found it. This section is
required. ``Everything worked first try'' scores zero.
\end{itemize}

\section*{7. Grading (100 points)}

\begin{center}
\begin{tabular}{lr}
\toprule
Pre-lab, submitted at the start of the session & 5 \\
Design compiles cleanly, warnings explained & 25 \\
Testbench and test vectors, coverage justified & 20 \\
Simulation results match the specification & 20 \\
Report chapter quality & 15 \\
In-lab demonstration and explanation & 15 \\
\midrule
Total & 100 \\
\bottomrule
\end{tabular}
\end{center}

Late work is accepted for one week at $-20$ points.

\section*{8. Things that will go wrong}

\begin{itemize}[leftmargin=1.6em,itemsep=1pt]
\item \textbf{Incomplete sensitivity list.} Simulates wrong, synthesises right.
\item \textbf{\code{=} versus \code{<=}.} \code{if En = '1'} is a comparison;
\code{Q <= D} is a signal assignment. Same symbols, different jobs.
\item \textbf{Expecting edge behaviour.} A latch is level-sensitive. It has no
clock and no edge. That arrives next week.
\item \textbf{Clearing a vector with \code{'0'}.} Use \code{(others => '0')}.
\item \textbf{Project path.} Short, local, no spaces. Not in OneDrive.
\end{itemize}

\vspace{10pt}
\hr
{\small This latch is the storage primitive for the entire semester. Lab 2 builds a
flip-flop from two of them; Lab 6 builds 1024. Keep your files.}

\end{document}
